% API reference for the laupy Python package.
% Compile with:  tectonic laupy_api_reference.tex
\documentclass[11pt,a4paper]{article}

\usepackage{fontspec}
\usepackage[margin=2.3cm]{geometry}
\usepackage{xcolor}
\usepackage[most]{tcolorbox}
\usepackage{listings}
\usepackage{booktabs,tabularx,longtable,array}
\usepackage{enumitem}
\usepackage{titlesec}
\usepackage{fancyhdr}
\usepackage[hidelinks]{hyperref}

\definecolor{accent}{HTML}{1F5A96}
\definecolor{termbg}{HTML}{F4F6F8}
\definecolor{termframe}{HTML}{B8C2CC}
\definecolor{notebg}{HTML}{EAF2FB}
\definecolor{warnbg}{HTML}{FDF1E4}
\definecolor{warnframe}{HTML}{D9822B}
\definecolor{sigbg}{HTML}{E8EEF5}
\definecolor{kw}{HTML}{1F5A96}
\definecolor{cm}{HTML}{6A737D}
\definecolor{st}{HTML}{8A3B12}

\hypersetup{colorlinks=true,linkcolor=accent,urlcolor=accent}
\setlist{itemsep=2pt,topsep=4pt}
\setlength{\parindent}{0pt}
\setlength{\parskip}{5pt}
\setlength{\emergencystretch}{3em}

\titleformat{\section}{\Large\bfseries\sffamily\color{accent}}{\thesection}{0.8em}{}
\titleformat{\subsection}{\large\bfseries\sffamily}{\thesubsection}{0.8em}{}
\titleformat{\subsubsection}{\normalsize\bfseries\sffamily}{\thesubsubsection}{0.8em}{}

\pagestyle{fancy}
\fancyhf{}
\fancyhead[L]{\small\sffamily laupy --- API Reference}
\fancyhead[R]{\small\sffamily\nouppercase{\leftmark}}
\fancyfoot[C]{\small\thepage}
\renewcommand{\headrulewidth}{0.4pt}

\lstdefinestyle{py}{
  language=Python,
  basicstyle=\ttfamily\small,
  keywordstyle=\color{kw}\bfseries,
  commentstyle=\color{cm}\itshape,
  stringstyle=\color{st},
  breaklines=true,
  postbreak=\mbox{\textcolor{gray}{$\hookrightarrow$}\space},
  columns=fullflexible,
  keepspaces=true,
  showstringspaces=false,
  upquote=true,
}
\lstset{style=py}

\newtcblisting{code}[1][]{
  listing only, colback=termbg, colframe=termframe, boxrule=0.4pt, arc=2pt,
  left=4pt, right=4pt, top=2pt, bottom=2pt, breakable,
  listing options={style=py}, #1
}
% no optional argument: looking for one would read a first line starting with # as TeX
\newtcblisting{terminal}{
  listing only, colback=termbg, colframe=termframe, boxrule=0.4pt, arc=2pt,
  left=4pt, right=4pt, top=2pt, bottom=2pt, breakable,
  listing options={language={},basicstyle=\ttfamily\small,commentstyle={},stringstyle={},keywordstyle={}}
}
\newtcolorbox{note}[1][Note]{
  colback=notebg, colframe=accent, boxrule=0.4pt, arc=2pt, breakable,
  fonttitle=\bfseries\sffamily\small, title={#1}, left=5pt, right=5pt
}
\newtcolorbox{warning}[1][Caution]{
  colback=warnbg, colframe=warnframe, boxrule=0.4pt, arc=2pt, breakable,
  fonttitle=\bfseries\sffamily\small, title={#1}, left=5pt, right=5pt
}

% API entry: signature box followed by the description
\newcommand{\api}[1]{%
  \par\medskip
  \begin{tcolorbox}[colback=sigbg, colframe=sigbg, boxrule=0pt, arc=1pt, left=4pt, right=4pt, top=2pt, bottom=2pt, breakable]
  \raggedright\ttfamily\small #1
  \end{tcolorbox}\nopagebreak\vspace{-4pt}%
}
% Allow line breaks after underscores and dots in long identifiers
\DeclareRobustCommand{\_}{\textunderscore\allowbreak}
\newcommand{\cmd}[1]{\texttt{#1}}
\newcommand{\args}{\par\textsf{\small\bfseries Arguments}\par}
\newcommand{\returns}{\par\textsf{\small\bfseries Returns}\quad}
\newcommand{\raises}{\par\textsf{\small\bfseries Raises}\quad}

\begin{document}

\begin{titlepage}
  \centering
  \vspace*{3cm}
  {\sffamily\Huge\bfseries\color{accent} laupy\par}
  \vspace{0.6cm}
  {\sffamily\LARGE API Reference\par}
  \vspace{1.2cm}
  {\large Python interface for SLURM pipelines of processing units\par}
  \vspace{2cm}
  {\large version 1.2.5\par}
  \vfill
  {\small Companion to the \emph{laupy \& lautools-browser User Manual}\par}
  {\small Licensed under the GNU General Public License v3.0\par}
  {\small Compiled \today\par}
\end{titlepage}

\tableofcontents
\clearpage

% ======================================================================
\section{Introduction}
% ======================================================================

laupy submits SLURM batch scripts for many \emph{processing units} (one sample each) and
records each submitted job as an \emph{entry} in the processing unit's
\cmd{pipeline/dag.json}. The command line tools \cmd{submitslurm} and \cmd{pipeline} and
lautools-browser are built on the API described here; your own scripts can use the same API.

\subsection{Modules}

\begin{center}
\begin{tabularx}{\linewidth}{@{}l>{\raggedright\arraybackslash}X@{}}
\toprule
Module & Purpose \\
\midrule
\cmd{laupy.pipeline\_entries} & Lifecycle of entries after submission: \cmd{PipelineEntriesStore} (load many processing units, one batched SLURM refresh, merging save), selection, retire, cancel. \textbf{Start here.} \\
\cmd{laupy.flow} & Files and scheduling: load/store \cmd{dag.json}, list processing units, refresh entries, build the DAG, submit and resubmit jobs. \\
\cmd{laupy.slurm} & Thin, batched wrappers around \cmd{sacct}, \cmd{squeue}, \cmd{sbatch}, \cmd{scancel}; job state constants. \\
\cmd{laupy.cli} & Helpers shared by the command line tools: common options, directory tree resolution, job environment. \\
\cmd{laupy.graph\_utils} & Minimal directed graph with ancestors, cycle detection and topological sort. \\
\cmd{laupy.maxwell} & Maxwell node categories (GPU/CPU, excellent/strong/weak) from a node table and \cmd{sinfo}; cluster topology report. \\
\cmd{laupy.redis\_server} & Lifecycle of the temporary Redis server of the GPU pool: start, stop, status (used by \cmd{REDIS} and \cmd{GPU -{}-try-start-server}). \\
\cmd{laupy.util} & Process listing used by \cmd{kct-list}. \\
\bottomrule
\end{tabularx}
\end{center}

\subsection{Data model}

\begin{terminal}
<project root>/                       working directory of sbatch
|-- sbatch/                           SLURM scripts
`-- wd/                               working directory
    `-- <processing unit subdir>/     one sample
        `-- pipeline/
            |-- dag.json              PIPELINE_ENTRIES: JSON list of entries
            |-- exec.sh               commands submitted (written by submitslurm)
            `-- log/                  <script>-<node>-<slurmid>.out/.err
\end{terminal}

\begin{description}[style=nextline,leftmargin=1.5em]
  \item[\cmd{PIPELINE\_ENTRIES}] The list of entry dicts of \emph{one} processing unit, as stored
    in \cmd{dag.json}. Functions take it as their first argument and change it \emph{in place}.
    The upper-case name is a convention throughout laupy.
  \item[Entry] One submitted SLURM job, a \cmd{dict}; fields in Section~\ref{sec:fields}. It is
    identified by \cmd{job\_id}, a 32-character hex UUID assigned by laupy, and carries the
    \cmd{slurm\_id} returned by \cmd{sbatch}.
  \item[DAG node] Every entry belongs to a node \cmd{dag\_node\_name}. Numeric nodes
    \cmd{"0"}, \cmd{"1"}, \ldots{} form a chain; named nodes depend on the nodes in
    \cmd{dag\_node\_deps}, or on the root when that list is empty; node \cmd{"-1"} is outside the
    DAG (Section~\ref{sec:dag}).
  \item[\cmd{slurm\_info}] The last known SLURM data of an entry: \cmd{State}, \cmd{JobName},
    \cmd{Elapsed}, \cmd{NodeList}, \cmd{StdOut}, \ldots{} (Section~\ref{sec:slurminfo}). It is
    updated by a \emph{refresh}; once the state is final (\cmd{TERMINAL\_STATES}), it is not
    queried again by default.
  \item[Retired] \cmd{entry["retired"] = True}: kept for history, ignored by scheduling, status and
    dependency resolution.
\end{description}

\subsection{The load $\rightarrow$ modify $\rightarrow$ store pattern}

All changes follow the same pattern: load the entries, refresh them from SLURM if states matter,
change the lists in place, store them.

\begin{code}
from laupy.pipeline_entries import PipelineEntriesStore

store = PipelineEntriesStore(subdirs)  # load_pipeline_entries() of every subdir
store.refresh()                        # one batched sacct/squeue query for all of them
for subdir, PIPELINE_ENTRIES in store.items():
    ...                                # read or change PIPELINE_ENTRIES in place
store.store()                          # write the changed subdirs, merged with the file on disk
\end{code}

For a single processing unit the same works with the functions of \cmd{laupy.flow}:
\cmd{load\_pipeline\_entries()}, \cmd{update\_pipeline\_entries()},
\cmd{store\_pipeline\_entries()}.

\begin{note}[Concurrent processes]
\cmd{dag.json} is replaced atomically, so readers never see a half-written file.
\cmd{PipelineEntriesStore.store()} also \emph{merges}: it re-reads the file, applies only the fields
you changed and appends your new entries, so entries appended meanwhile by another process (e.g.\
\cmd{submitslurm} while the browser refreshes) are kept. If two processes change the same field of
the same entry, the last write wins. \cmd{store\_pipeline\_entries()} does \emph{not} merge; it
overwrites the file with the list you pass.
\end{note}

% ======================================================================
\section{Examples}
% ======================================================================

All examples run in the project root and were tested with laupy 1.2.5.

\subsection{Status of all processing units}

\begin{code}
from laupy.flow import list_processing_unit_subdirs
from laupy.pipeline_entries import PipelineEntriesStore, entry_slurm_id

subdirs = list_processing_unit_subdirs("wd", require_pipeline=True)
store = PipelineEntriesStore(subdirs, topological_sort=True)
queried = store.refresh()              # one batched sacct/squeue call for all subdirs
for subdir, PIPELINE_ENTRIES in store.items():
    for entry in PIPELINE_ENTRIES:
        if entry.get("retired"):
            continue
        state = entry.get("slurm_info", {}).get("State", "UNKNOWN")
        print(entry["processing_unit_name"], entry["dag_node_name"], entry_slurm_id(entry), state)
store.store()                          # caches final states, writes only changed subdirs
print("queried", queried, "entries")
\end{code}

\begin{terminal}
syn001_ti_12w 1 24170101 COMPLETED
syn001_ti_12w 2 24170104 RUNNING
syn002_mg_4w 1 24170102 RUNNING
syn002_mg_4w 2 24170105 PENDING
syn003_peek_8w 1 24170103 FAILED
syn003_peek_8w 2 24170106 PENDING
queried 6 entries
\end{terminal}

\subsection{Retire failed entries and cancel chosen jobs}

\begin{code}
from laupy.flow import list_processing_unit_subdirs
from laupy.pipeline_entries import (PipelineEntriesStore, RETIRABLE_STATES, CANCELLABLE_STATES,
                                    select_pipeline_entries, retire_pipeline_entries, cancel_selected)

store = PipelineEntriesStore(list_processing_unit_subdirs("wd", require_pipeline=True))
store.refresh()
# Retire the failed entries of step 1
for subdir, PIPELINE_ENTRIES in store.items():
    failed = select_pipeline_entries(PIPELINE_ENTRIES, RETIRABLE_STATES, steps=[1])
    for entry in retire_pipeline_entries(PIPELINE_ENTRIES, [e["job_id"] for e in failed],
                                         reason="bad input data"):
        print("retired", entry["slurm_id"])
store.store()

# Cancel two jobs given by SLURM ID, wherever they are
chosen = select_pipeline_entries(store.all_entries(), CANCELLABLE_STATES,
                                 job_ids=["24170102", "24170105"])
for subdir, results in cancel_selected(chosen).items():
    for r in results:
        print(r["slurm_id"], r["state"], "->", r["error"] or "cancelled")
\end{code}

\begin{terminal}
retired 24170103
24170102 RUNNING -> cancelled
24170105 PENDING -> cancelled
\end{terminal}

\subsection{Submitting jobs}

\cmd{submitslurm} is the reference for submissions; this is its core, reduced to one node:

\begin{code}
import os
from laupy.cli import SubmissionEnvironment
from laupy.flow import list_processing_unit_subdirs, SchedulingError, schedule_slurm_job
from laupy.pipeline_entries import PipelineEntriesStore

root = os.getcwd()                                   # project root with sbatch/ and wd/
store = PipelineEntriesStore(list_processing_unit_subdirs("wd"))
store.refresh(limit_to_ancestors_including_node="1") # states of node 1 and its ancestors
env = SubmissionEnvironment(extra={"OMP_NUM_THREADS": "8"})

for subdir in store.subdirs():
    name = os.path.basename(subdir)
    script = "01createZarrStructure.sbatch"
    job_name = f"01createZarrStructure_{name}"
    log_dir = os.path.join(subdir, "pipeline", "log")
    os.makedirs(log_dir, exist_ok=True)
    cmd = [os.path.join(root, "sbatch", script), subdir, "--compression", "htj2k"]
    slurm_cmd = ["sbatch", "--partition=allcpu,allgpu", f"--job-name={job_name}",
                 env.export_option(env.job_vars(script, job_name)),
                 f"--output={log_dir}/01createZarrStructure-%N-%j.out",
                 f"--error={log_dir}/01createZarrStructure-%N-%j.err"] + cmd
    PIPELINE_ENTRIES = store[subdir]
    try:
        entry = schedule_slurm_job(PIPELINE_ENTRIES, "1", script, job_name, subdir,
                                   slurm_cmd, cmd, env=env.global_env, raise_errors=True)
    except SchedulingError as e:
        print(f"{name}: not submitted, {e}")
        continue
    PIPELINE_ENTRIES.append(entry)          # schedule_slurm_job does not append
    print(f"{name}: submitted {entry['slurm_id']}")
store.store()
\end{code}

\cmd{schedule\_slurm\_job()} adds \cmd{-{}-dependency=afterok:\ldots} for the unfinished jobs of
all ancestor nodes. It needs their current states, hence the refresh limited to the node and its
ancestors before submitting.

\subsection{One processing unit, resubmission and the DAG}

\begin{code}
from laupy.flow import (load_pipeline_entries, store_pipeline_entries, update_pipeline_entries,
                        constructDAG, resubmit_slurm_job)
from laupy.graph_utils import ancestors

PIPELINE_ENTRIES = load_pipeline_entries("wd/syn003_peek_8w")
update_pipeline_entries(PIPELINE_ENTRIES, update_retired=True)
print(sorted(ancestors(constructDAG(PIPELINE_ENTRIES), "2")))   # ['0', '1', 'dag_root']

failed = next(e for e in PIPELINE_ENTRIES if e["slurm_id"] == 24170103)
new_entry = resubmit_slurm_job(PIPELINE_ENTRIES, failed)          # old entry retired: "Resubmitted"
PIPELINE_ENTRIES.append(new_entry)
store_pipeline_entries("wd/syn003_peek_8w", PIPELINE_ENTRIES)
\end{code}

\subsection{Querying SLURM directly}

\begin{code}
from laupy import slurm

info = slurm.slurm_info([24170101, 24170104])      # list in, list out
print({i["JobID"]: i["State"] for i in info})       # {'24170101': 'COMPLETED', '24170104': 'RUNNING'}
print(slurm.QUERY_COUNTS)                           # {'sacct': 1, 'squeue': 1}
error = slurm.scancel(24170104)                     # None on success
\end{code}

% ======================================================================
\section{\texorpdfstring{\cmd{laupy.pipeline\_entries}}{laupy.pipeline\_entries}}
% ======================================================================

Lifecycle of entries after submission:
\emph{submitted} $\rightarrow$ \emph{refreshed} $\rightarrow$ \emph{cancelled / finished}
$\rightarrow$ \emph{retired}.
Shared by \cmd{pipeline}, \cmd{submitslurm} and lautools-browser.

\subsection{Constants}

\begin{center}
\begin{tabularx}{\linewidth}{@{}l>{\raggedright\arraybackslash}X@{}}
\toprule
Name & Value \\
\midrule
\cmd{RETIRABLE\_STATES} & \cmd{FAILED\_STATES}: FAILED, CANCELLED, TIMEOUT. \\
\cmd{CANCELLABLE\_STATES} & \cmd{ACTIVE\_STATES}: PENDING, RUNNING, REQUEUED, COMPLETING. \\
\bottomrule
\end{tabularx}
\end{center}

\subsection{Class \texorpdfstring{\cmd{PipelineEntriesStore}}{PipelineEntriesStore}}

\api{PipelineEntriesStore(subdirs=(), topological\_sort=False)}
The \cmd{PIPELINE\_ENTRIES} of several processing unit subdirs, with one batched SLURM refresh
and a merging \cmd{store()}. Subdirs are keyed by their resolved absolute path; every method
accepting a subdir also accepts a relative path or \cmd{Path}.
\args
\begin{description}[leftmargin=2em,font=\ttfamily\mdseries]
  \item[subdirs] Processing unit subdirs to load immediately.
  \item[topological\_sort] Order each list by the topological order of the DAG nodes, then by
    submission time (as \cmd{pipeline status} shows them).
\end{description}
On loading, each entry gets \cmd{processing\_unit\_subdir} and \cmd{processing\_unit\_name} if
missing; this alone does not mark the subdir as changed.

\api{load(subdir) -> list}
(Re)load one subdir from disk, discarding unsaved changes. Returns its list.

\api{subdirs() -> list[str]\quad items() -> Iterator[(str, list)]\quad all\_entries() -> list[dict]}
The loaded subdirs (absolute paths); pairs of subdir and \cmd{PIPELINE\_ENTRIES}; all entries of
all subdirs in one flat list. Entries are the stored objects, so changes made through
\cmd{all\_entries()} are saved by \cmd{store()} too.

\api{entries(subdir) -> list\quad store[subdir]}
The loaded list of one subdir. Change it in place (including \cmd{append()}), then \cmd{store()}.

\api{find(job\_id) -> dict | None}
The entry with this laupy \cmd{job\_id} in any subdir.

\api{refresh(subdirs=None, select=None, **selection) -> int}
Refresh \cmd{slurm\_info} of entries from all (or the given) subdirs with one batched query.
\args
\begin{description}[leftmargin=2em,font=\ttfamily\mdseries]
  \item[subdirs] Limit the refresh to these subdirs.
  \item[select] \cmd{select(subdir, PIPELINE\_ENTRIES) -> list}: choose the entries to query
    yourself (e.g.\ by \cmd{job\_id}).
  \item[**selection] Otherwise the entries are chosen by
    \cmd{laupy.flow.pipeline\_entries\_to\_refresh(PIPELINE\_ENTRIES, **selection)}, e.g.\
    \cmd{update\_retired=True}, \cmd{filter\_terminal\_states=False},
    \cmd{limit\_to\_ancestors\_including\_node="3"}.
\end{description}
\returns The number of entries queried.

\api{modified\_subdirs() -> list[str]}
Subdirs whose entries differ from what was loaded (field changes or new entries).

\api{store(subdirs=None) -> list[str]}
Write the changed subdirs (all, or only the given ones), merged with the current file on disk:
fields you changed are applied, new entries are appended, everything else on disk is kept.
Afterwards the lists hold the merged state. Entries you removed from a list are \emph{not} removed
from the file; use \cmd{laupy.flow.clean\_pipeline\_entries()} to forget entries.
\returns The subdirs written.
\raises \cmd{OSError} when a file cannot be written (e.g.\ read-only project).

\subsection{Functions}

\api{select\_pipeline\_entries(PIPELINE\_ENTRIES, states, job\_ids=None, steps=None) -> list}
Non-retired entries whose \cmd{slurm\_info["State"]} is in \cmd{states}.
\args
\begin{description}[leftmargin=2em,font=\ttfamily\mdseries]
  \item[PIPELINE\_ENTRIES] Any list of entries, also \cmd{store.all\_entries()}.
  \item[states] Iterable of state names, e.g.\ \cmd{RETIRABLE\_STATES}.
  \item[job\_ids] Limit to entries whose laupy \cmd{job\_id} \emph{or} SLURM ID is in this list;
    both kinds may be mixed, compared as strings.
  \item[steps] Limit to these step numbers (used only when \cmd{job\_ids} is not given).
\end{description}
Refresh first: the selection uses the stored states.

\api{entry\_matches\_id(entry, ids) -> bool}
True if the entry's \cmd{job\_id} or SLURM ID is in \cmd{ids} (a set of strings).

\api{entry\_slurm\_id(entry)}
The entry's \cmd{slurm\_id}, or the \cmd{JobID} reported by SLURM if missing.

\api{retire\_pipeline\_entries(PIPELINE\_ENTRIES, job\_ids, reason=None, retired\_by=None) -> list}
Retire the not yet retired entries with these laupy \cmd{job\_id}s. Sets \cmd{retired},
\cmd{retire\_timestamp} and, if given, \cmd{retired\_reason} and \cmd{retired\_by} (by convention
the \cmd{job\_id} of the entry that replaced them). Does not save.
\returns The entries retired by this call.

\api{cancel\_pipeline\_entries(PIPELINE\_ENTRIES, job\_ids, refresh=True) -> list[dict]}
\cmd{scancel} the entries with these laupy \cmd{job\_id}s and set their state to CANCELLED. Entries
are not retired. With \cmd{refresh=True} their states are queried first (one batched query); jobs no
longer active are reported, not cancelled. Each job is cancelled separately, so the outcome is
known per job. Does not save.
\returns One dict per entry: \cmd{entry}, \cmd{slurm\_id}, \cmd{state} (before cancelling),
\cmd{error} (message or \cmd{None}).

\api{retire\_selected(selection, reason=None) -> dict[str, list]}
Retire entries chosen elsewhere (e.g.\ in a dialog): the affected subdirs are loaded fresh, changed
and saved.
\cmd{selection} is \cmd{\{subdir: [job\_id, \ldots]\}} or a list of entries carrying
\cmd{processing\_unit\_subdir} (e.g.\ from \cmd{store.all\_entries()}).
\returns \cmd{\{subdir: retired entries\}}.

\api{cancel\_selected(selection) -> dict[str, list]}
Cancel chosen entries of any number of subdirs with one SLURM query, then save (including the
refreshed states). \cmd{selection} as for \cmd{retire\_selected()}.
\returns \cmd{\{subdir: results of cancel\_pipeline\_entries()\}}.

% ======================================================================
\section{\texorpdfstring{\cmd{laupy.flow}}{laupy.flow}}
% ======================================================================

\subsection{Files and directories}

\api{pipeline\_entries\_file(processing\_unit\_subdir) -> str}
Path of \cmd{<subdir>/pipeline/dag.json}.

\api{load\_pipeline\_entries(processing\_unit\_subdir, topological\_sort=False) -> list}
The entries of one processing unit; \cmd{[]} if there is no file. Field names written by laupy
before 1.2.5 (\cmd{execution\_unit\_dir}, \cmd{execution\_unit\_name}) are converted to
\cmd{processing\_unit\_subdir}, \cmd{processing\_unit\_name}.

\api{store\_pipeline\_entries(processing\_unit\_subdir, PIPELINE\_ENTRIES) -> None}
Write the list, replacing the file atomically (temporary file and \cmd{os.replace}). Creates
\cmd{pipeline/} if needed. Overwrites; see \cmd{PipelineEntriesStore.store()} for merging.

\api{sort\_pipeline\_entries\_topologically(PIPELINE\_ENTRIES) -> list}
A new list ordered by the topological order of the DAG nodes, then by \cmd{timestamp}.

\api{clean\_pipeline\_entries(processing\_unit\_subdir) -> None}
Forget all entries (empty \cmd{dag.json}); logs are kept. As \cmd{pipeline clean}.

\api{list\_processing\_unit\_subdirs(working\_dir, pattern=None, require\_pipeline=False) -> list[Path]}
The immediate subdirectories of a working directory, sorted by name.
\args
\begin{description}[leftmargin=2em,font=\ttfamily\mdseries]
  \item[pattern] Case-insensitive substring the name must contain (as \cmd{-p}).
  \item[require\_pipeline] Only subdirectories that already have \cmd{pipeline/}.
\end{description}

\api{remove\_pipeline\_dir(pipeline\_dir) -> None}
Delete the files of \cmd{pipeline/} and \cmd{pipeline/log/} and both directories. Other
subdirectories of \cmd{pipeline/} are left untouched. As \cmd{pipeline delete}.

\subsection{Refreshing entries}

\api{pipeline\_entries\_to\_refresh(PIPELINE\_ENTRIES, limit\_to\_ancestors\_including\_node="-1", update\_retired=False, update\_negative\_step=False, filter\_terminal\_states=True, node\_dependencies=None) -> list}
Select the entries whose SLURM state should be queried.
\args
\begin{description}[leftmargin=2em,font=\ttfamily\mdseries]
  \item[limit\_to\_ancestors\_including\_node] Only entries of this node and its ancestors;
    \cmd{"-1"}: all entries. Ancestors are computed on the full DAG.
  \item[update\_retired] Include retired entries.
  \item[update\_negative\_step] Include entries of step $-1$.
  \item[filter\_terminal\_states] Skip entries whose stored state is final (\cmd{TERMINAL\_STATES}).
  \item[node\_dependencies] Dependencies of a node not yet in the entries (a new named node), so
    that its ancestors can be determined before its first submission.
\end{description}

\api{refresh\_pipeline\_entries(entries) -> None}
Query SLURM for the given entries, from any number of processing units, in one batched query and
store the result in their \cmd{slurm\_info}. Entries SLURM no longer knows keep their stored info.

\api{update\_pipeline\_entries(PIPELINE\_ENTRIES, **selection) -> None}
\cmd{refresh\_pipeline\_entries(pipeline\_entries\_to\_refresh(PIPELINE\_ENTRIES, **selection))}
for one processing unit. For many units use \cmd{PipelineEntriesStore.refresh()}.

\subsection{The DAG}\label{sec:dag}

\api{constructDAG(PIPELINE\_ENTRIES) -> DirectedGraph}
The DAG of the nodes referenced by the entries:
\begin{itemize}
  \item a root node \cmd{"dag\_root"} with the edge \cmd{dag\_root} $\rightarrow$ \cmd{"0"};
  \item numeric nodes form a chain \cmd{"0"} $\rightarrow$ \cmd{"1"} $\rightarrow$ \ldots{} up to the
    largest number used; they cannot have explicit dependencies;
  \item named nodes get edges from each node in their \cmd{dag\_node\_deps} (unknown dependencies are
    dropped with a warning), or from \cmd{dag\_root} when the list is empty;
  \item negative nodes (\cmd{"-1"}) are not part of the graph.
\end{itemize}
Entries without \cmd{dag\_node\_name} get it from \cmd{step} (normalised to a string, in place).
\raises \cmd{ValueError} if the graph contains a cycle.

\api{string\_is\_integer(s) -> bool}
True if \cmd{int(s)} succeeds.

\subsection{Scheduling}

\api{schedule\_slurm\_job(PIPELINE\_ENTRIES, dag\_node\_or\_step\_name, script\_name, job\_name, processing\_unit\_subdir, slurm\_cmd, cmd, node\_dependencies=None, env=None, raise\_errors=False, laupy\_python=None) -> dict | None}
Submit one job for a DAG node and return its new entry. The caller appends it to
\cmd{PIPELINE\_ENTRIES} and stores.
\args
\begin{description}[leftmargin=2em,font=\ttfamily\mdseries]
  \item[dag\_node\_or\_step\_name] \cmd{"3"}, \cmd{"reco"} or \cmd{"-1"} (no dependencies).
  \item[script\_name, job\_name] Recorded in the entry.
  \item[processing\_unit\_subdir] Absolute path of the unit; \cmd{sbatch} runs in its
    grandparent (the project root).
  \item[slurm\_cmd] Full \cmd{sbatch} argument list, ending with the script command; stored for
    resubmission.
  \item[cmd] The script command (script and arguments); stored, space-joined, as \cmd{command}.
  \item[node\_dependencies] Parents of a named node.
  \item[env] Environment for \cmd{sbatch} (e.g.\ \cmd{SubmissionEnvironment.global\_env}).
  \item[raise\_errors] Raise \cmd{SchedulingError} instead of logging and returning \cmd{None}.
  \item[laupy\_python] Python environment of the job, recorded as \cmd{laupy\_python}; \cmd{None}
    records nothing. It is not added to \cmd{slurm\_cmd}: the caller puts \cmd{LAUPY\_PYTHON} into its
    \cmd{-{}-export}.
\end{description}
For non-negative nodes, all non-retired, not completed entries of ancestor nodes are checked with
\cmd{slurm.get\_active\_slurmids()}: active ones become \cmd{-{}-dependency=afterok:\ldots}, recorded in
\cmd{dependency\_slurm\_ids}; a failed one stops the submission. The states used are the stored
ones, so refresh the ancestors first.
\raises \cmd{SchedulingError} with a one-line reason (only with \cmd{raise\_errors=True}).

\api{resubmit\_slurm\_job(PIPELINE\_ENTRIES, entry\_id, env=None) -> dict | None}
Submit an entry again with its stored \cmd{slurm\_command}. If its job is still active, it is
cancelled first. The old entry is retired (\cmd{retired\_reason="Resubmitted"},
\cmd{retired\_by=<new job\_id>}) only after the new submission succeeded.
\cmd{entry\_id} is the entry dict, its \cmd{slurm\_id} (\cmd{int} or string of digits) or its
\cmd{job\_id}; surrounding spaces are ignored and \cmd{bool} is rejected. Entries without a
\cmd{slurm\_id} never match a string ID. The caller appends the new entry and stores.
The new entry keeps \cmd{laupy\_python}. For an entry without it, \cmd{LAUPY\_PYTHON} is removed from
\cmd{env} (or from \cmd{os.environ}), so the job gets the default environment as originally.
\returns The new entry, or \cmd{None} after logging the reason when the entry is not found, is
already retired (resubmitting it would start a second job for the same node), or the submission
fails.

\api{class SchedulingError(RuntimeError)}
A job could not be scheduled: failed dependency, \cmd{sbatch} failure or invalid DAG.

% ======================================================================
\section{\texorpdfstring{\cmd{laupy.slurm}}{laupy.slurm}}
% ======================================================================

\subsection{Constants}

\begin{center}
\begin{tabularx}{\linewidth}{@{}l>{\raggedright\arraybackslash}X@{}}
\toprule
Name & Value \\
\midrule
\cmd{ACTIVE\_STATES} & PENDING, RUNNING, REQUEUED, COMPLETING \\
\cmd{FAILED\_STATES} & FAILED, CANCELLED, TIMEOUT \\
\cmd{TRANSIENT\_STATES} & COMPLETING: normally seconds; cancel only if stuck. \\
\cmd{TERMINAL\_STATES} & COMPLETED, FAILED, CANCELLED, TIMEOUT, OUT\_OF\_MEMORY, BOOT\_FAIL, DEADLINE: final, not queried again. \\
\cmd{QUERY\_CHUNK\_SIZE} & 500: IDs per \cmd{sacct}/\cmd{squeue} call. \\
\cmd{QUERY\_COUNTS} & \cmd{\{"sacct": n, "squeue": n\}}: calls made by this process, for diagnostics. \\
\bottomrule
\end{tabularx}
\end{center}

\subsection{Job information}\label{sec:slurminfo}

\api{slurm\_info(slurm\_ids) -> dict | list[dict]}
Job information from \cmd{sacct}, completed by \cmd{squeue} for pending and running jobs. One call
of each tool per \cmd{QUERY\_CHUNK\_SIZE} unique IDs. A single ID returns a dict, a list returns a
list in the same order. IDs that are not SLURM IDs (\cmd{\textbackslash d+} or
\cmd{\textbackslash d+\_\textbackslash d+}) are not sent, because one invalid ID would make the
whole call fail; unknown jobs get \cmd{State="UNKNOWN"}.

Keys from \cmd{sacct}: \cmd{JobID}, \cmd{JobName}, \cmd{State}, \cmd{ExitCode}, \cmd{Partition},
\cmd{User}, \cmd{Elapsed}, \cmd{Timelimit}, \cmd{NNodes}, \cmd{NodeList}, \cmd{Reason}, \cmd{StdOut},
\cmd{StdErr} (with \cmd{\%j}, \cmd{\%N} replaced). \cmd{squeue} updates \cmd{NodeList}, \cmd{State},
\cmd{Reason} and \cmd{Elapsed} of queued jobs. States are normalised
(e.g.\ \emph{CANCELLED by 1234} $\rightarrow$ CANCELLED).

\api{slurm\_sacct\_info(slurm\_ids)\quad slurm\_squeue\_info(slurm\_ids)}
The two halves of \cmd{slurm\_info()}. \cmd{slurm\_squeue\_info()} reports jobs no longer queued as
UNKNOWN.

\api{normalize\_state(state) -> str}
First word of a SLURM state without a trailing \cmd{+}: \emph{CANCELLED by 1234} $\rightarrow$ CANCELLED.

\api{get\_active\_slurmids(dependency\_entries, raise\_on\_fail=False) -> list}
The \cmd{slurm\_id}s of the entries whose stored state is in \cmd{ACTIVE\_STATES}. Entries in
\cmd{FAILED\_STATES} print a warning, or raise \cmd{RuntimeError} with
\cmd{raise\_on\_fail=True}. Other states count as finished.

\api{list\_user\_jobs(user=None) -> list[dict]}
All queued and running jobs of a user (default: the current one) from \cmd{squeue}. Keys:
\cmd{JobID}, \cmd{JobName}, \cmd{State}, \cmd{Partition}, \cmd{Elapsed}, \cmd{Timelimit},
\cmd{Reason}, \cmd{SubmitTime}, \cmd{WorkDir}.
\raises \cmd{subprocess.CalledProcessError} or \cmd{OSError} when \cmd{squeue} fails.

\api{own\_allocation\_job\_ids() -> set\quad is\_own\_allocation(slurm\_id, own\_ids=None) -> bool}
The SLURM job the current process runs in (from \cmd{SLURM\_JOB\_ID}/\cmd{SLURM\_JOBID}; empty
outside SLURM), and whether an ID (also an array task) is that job. Used to protect the caller's
own allocation from \cmd{scancel}.

\subsection{Submitting and cancelling}

\api{run\_slurm\_command(SLURM\_CMD, cwd=None, env=None) -> int | None}
Run an \cmd{sbatch} command and return the job ID parsed from \emph{Submitted batch job N}.
Returns \cmd{None} and logs the command, working directory, stdout and stderr on failure, or when
the output contains no job ID (the outcome is then uncertain; check SLURM before retrying).

\api{scancel(slurm\_ids) -> str | None}
One \cmd{scancel} call for one or more IDs. \cmd{None} on success, otherwise the error message.
\cmd{scancel} still cancels the valid IDs when some fail (e.g.\ already finished), so pass a single
ID when the outcome per job matters.

% ======================================================================
\section{\texorpdfstring{\cmd{laupy.cli}}{laupy.cli}}
% ======================================================================

Helpers shared by \cmd{pipeline} and \cmd{submitslurm}; useful for your own tools with the same
options.

\api{add\_common\_arguments(parser, suppress\_defaults=False, dry\_run\_help=...)}
Add \cmd{-d/-{}-root-dir}, \cmd{-w/-{}-working-dir} (comma-separated or repeated),
\cmd{-{}-slurm-dir}, \cmd{-p/-{}-pattern}, \cmd{-{}-dry-run}, \cmd{-{}-verbose}. With
\cmd{suppress\_defaults=True} (for a subcommand parent parser), the options only override values
given before the subcommand when given explicitly.

\api{add\_node\_source\_arguments(parser, offline=False)\quad apply\_node\_source\_arguments(parser, args)}
Add \cmd{-{}-list FILE}, \cmd{-{}-no-list}, \cmd{-{}-feature-provider} (and \cmd{-{}-offline}) choosing where
\cmd{laupy.maxwell} takes the nodes from, and apply them with \cmd{maxwell.set\_node\_table()}
(\cmd{parser.error} on conflicts or a missing file). Applying does not query SLURM.
\cmd{print\_node\_notes(file=None)} prints and clears \cmd{maxwell.NODE\_NOTES}.

\api{resolve\_directory\_tree(ARG) -> DirectoryTree}
Change to the root directory and resolve the working directories, the SLURM script directory and
the processing unit subdirs selected by \cmd{-w}/\cmd{-p}. Prints an error and exits when a
directory is missing or nothing matches; prints nothing otherwise.

\api{@dataclass DirectoryTree(root, sbatch\_abs, sbatch\_rel, working\_dirs\_rel, processing\_unit\_subdirs)}
\cmd{root: Path}; \cmd{sbatch\_abs: Path}, \cmd{sbatch\_rel: str}; \cmd{working\_dirs\_rel: list[str]};
\cmd{processing\_unit\_subdirs: list[ProcessingUnitSubdir]}.

\api{@dataclass ProcessingUnitSubdir(name, abs, rel)}
Directory name, resolved absolute path, path relative to the root (e.g.\ \cmd{wd/syn001}).

\api{@dataclass SubmissionEnvironment(global\_env=os.environ copy, extra=\{\}, vanilla=False)}
Environment of submitted jobs in two layers. \cmd{global\_env} is the environment \cmd{sbatch}
runs with; every job inherits it through \cmd{-{}-export=ALL}. \cmd{extra} holds additional
per-job variables (\cmd{submitslurm -{}-env}).
\begin{description}[leftmargin=2em,font=\ttfamily\mdseries]
  \item[job\_vars(script\_name, job\_name) -> dict] \cmd{PIPELINE\_SCRIPTNAME},
    \cmd{PIPELINE\_JOBNAME}, \cmd{LD\_PRELOAD=""}, then \cmd{extra} (which may override them).
  \item[export\_option(job\_vars) -> str] \cmd{"-{}-export=ALL,NAME=VALUE,\ldots"}, without \cmd{ALL} if \cmd{vanilla}.
  \item[describe\_job\_vars(job\_vars) -> str] Shell-quoted \cmd{NAME=VALUE} list for display.
  \item[global\_lines() -> list[str]] \cmd{CONDA\_PREFIX} and \cmd{PATH} of \cmd{global\_env} for
    display, see \cmd{environment\_lines()}.
\end{description}
With \cmd{vanilla=True} (\cmd{submitslurm -{}-vanilla} or \cmd{-{}-vanilla-strict}), \cmd{global\_env} is not propagated:
\cmd{export\_option()} lists the variables without \cmd{ALL}, and \cmd{job\_vars()} starts with the
\cmd{base} variables (field \cmd{base}, a dict; set it from \cmd{vanilla\_base\_variables()}) before
the usual ones. \cmd{export\_option()} is an instance method.

\api{vanilla\_base\_variables(root, sbatch\_dir, global\_env=None, job\_path=None, strict=False, commands=VANILLA\_PATH\_COMMANDS) -> dict}
\cmd{PATH} (\cmd{root/bin}, \cmd{root/script}, \cmd{sbatch\_dir}, the \cmd{job\_path} directories, the
directories of \cmd{vanilla\_kept\_directories(global\_env, commands)} unless \cmd{strict}
(\cmd{-{}-vanilla-strict}), \cmd{VANILLA\_SYSTEM\_PATH}), \cmd{HOME}, \cmd{USER},
\cmd{LOGNAME} (password database; \cmd{Path.home()}/\cmd{getpass.getuser()} if the UID has no entry).
\raises \cmd{ValueError} if a value contains a comma.

\api{vanilla\_kept\_directories(global\_env=None, commands=VANILLA\_PATH\_COMMANDS) -> (dict, list)}
For each command, the directory of its first occurrence on the \cmd{PATH} of \cmd{global\_env} (or
\cmd{os.environ}). Returns \cmd{\{directory: [commands]\}} in the order of that \cmd{PATH} and the
commands not found. \cmd{VANILLA\_PATH\_COMMANDS} is \cmd{("initEnv", "rawTiffStructureToZarr.py",
"kct-pb2d-krylov")}: the laupy \cmd{bin}, syscripts and the KCT binaries.

\api{parse\_job\_path(value) -> list[str]\quad job\_path\_variable(job\_path, global\_env=None, leading=()) -> str}
\cmd{argparse} type of \cmd{submitslurm -{}-job-path}: the colon-separated directories made absolute
(no comma). \cmd{job\_path\_variable()} gives the \cmd{PATH} of a non-vanilla job: \cmd{leading},
\cmd{job\_path}, then the inherited \cmd{PATH}, without duplicates; inherited directories with a comma
are left out with a warning (\cmd{ValueError} for a comma in \cmd{leading} or \cmd{job\_path}).

\api{workspace\_path\_directories(root, sbatch\_dir) -> list[str]\quad prepend\_path\_directories(leading, path\_value) -> str}
The directories \cmd{submitslurm} puts first on the \cmd{PATH} of every job:
\cmd{[root/bin, root/script, sbatch\_dir]}, labelled by \cmd{WORKSPACE\_PATH\_LABELS}.
\cmd{prepend\_path\_directories()} gives \cmd{path\_value} with \cmd{leading} first, each directory
once (used for \cmd{-{}-env PATH=}).

\api{environment\_lines(env, indent="    ") -> list[str]}
\cmd{CONDA\_PREFIX} and the \cmd{PATH} components, one per line; components below
\cmd{CONDA\_PREFIX} are abbreviated, duplicates and missing directories are flagged.

\api{check\_python\_environment(value) -> str\quad parse\_python\_environment(value) -> str}
Absolute path of a Conda/Mamba environment (a directory with executable \cmd{bin/python}, no comma) for
\cmd{LAUPY\_PYTHON}; the first raises \cmd{ValueError}, the second is the \cmd{argparse} type of
\cmd{submitslurm -{}-python}. \cmd{cli.LAUPY\_PYTHON} is the variable name.

\api{join\_dash\_values(argv, options=DASH\_VALUE\_OPTIONS) -> list[str]}
Rewrites \cmd{-a VALUE}, \cmd{-{}-slurmargs VALUE} and \cmd{-{}-env VALUE} as \cmd{-{}-long=VALUE} before
\cmd{argparse}, so a value starting with \cmd{-} (\cmd{-a "-{}-verbose"}) is not taken for an option.
Tokens after \cmd{-{}-} are left alone. \cmd{submitslurm} applies it to \cmd{sys.argv}.

\api{parse\_env\_assignment(value) -> (str, str)\quad parse\_comma\_separated(value) -> list[str]}
\cmd{argparse} types: \cmd{NAME=VALUE} (no commas in the value, since \cmd{-{}-export} separates
with commas) and \cmd{"a,b,c"} $\rightarrow$ \cmd{["a", "b", "c"]}.

\api{print\_arguments(ARG)}
Print the parsed arguments (\cmd{-{}-verbose}).

% ======================================================================
\section{\texorpdfstring{\cmd{laupy.submitslurm\_options}}{laupy.submitslurm\_options}}

The \cmd{submitslurm} option set as a component for programs that compose \cmd{submitslurm} calls
(the lautools submission context editor). \cmd{submitslurm} parses its own command line with
\cmd{build\_parser()}, so a new option is added there and every consumer sees it;
\cmd{tests/test\_submitslurm\_options.py} checks that \cmd{submitslurm} uses this parser.

\api{build\_parser() -> argparse.ArgumentParser}
The \cmd{submitslurm} parser, used by \cmd{submitslurm} itself.

\api{option\_specs(parser=None) -> list[OptionSpec]\quad context\_option\_specs(parser=None) -> list[OptionSpec]}
Every option (without help and positional), or only those of category \cmd{"context"}.
\cmd{OptionSpec} has \cmd{dest}, \cmd{option\_strings}, \cmd{name} (the long string), \cmd{kind}
(\cmd{flag}, \cmd{value}, \cmd{choice}, \cmd{multi\_choice}, \cmd{append}), \cmd{category}
(\cmd{context}, \cmd{call}: \cmd{CALL\_DESTS}, \cmd{action}: \cmd{ACTION\_DESTS}), \cmd{help},
\cmd{metavar}, \cmd{choices}, \cmd{default}.

\api{split\_arguments(argv, parser=None) -> list[ArgumentGroup]}
The call split into options with their values (\cmd{dest}, \cmd{tokens}) as \cmd{argparse} reads it;
handles \cmd{-{}-opt=value}, \cmd{-fVALUE} and dashed values of \cmd{-a}/\cmd{-{}-env}. Other tokens are
groups with \cmd{dest=None}.

\api{merge\_arguments(call, defaults=None, forced=frozenset(), extra=(), parser=None) -> list[str]}
\cmd{defaults} maps a \cmd{dest} to its tokens; they are used when the call lacks the option, or
always for \cmd{forced} dests (the call's own are dropped) and for \cmd{MERGED\_DESTS} (\cmd{env}: the
call's later \cmd{-{}-env} wins per name). Result: context tokens, \cmd{extra}, the call's options,
then its positional arguments after \cmd{-{}-} (so \cmd{-{}-skip-state FAILED} cannot take the script).

\api{parse\_merged(argv, parser=None) -> argparse.Namespace}
Parses like \cmd{submitslurm} (with \cmd{join\_dash\_values}); raises \cmd{ValueError} instead of exiting.

\api{preview\_environment(args, root, global\_env=None, script\_name="SCRIPT.sbatch", job\_name="SCRIPT\_UNIT") -> EnvironmentPreview}
The environment a job would get: \cmd{variables} (name $\rightarrow$ (value, source)), \cmd{path}
(list of (directory, source)), \cmd{export}, \cmd{inherits} (\cmd{-{}-export=ALL}) and \cmd{warnings}.
Computed with the same \cmd{laupy.cli} functions as \cmd{submitslurm}.

\section{\texorpdfstring{\cmd{laupy.graph\_utils}}{laupy.graph\_utils}}
% ======================================================================

A dependency-free replacement for the few \cmd{networkx} functions laupy uses.

\api{class DirectedGraph(edges=None)}
\cmd{add\_node(n)}, \cmd{add\_edge(u, v)}, \cmd{add\_edges\_from(edges)},
\cmd{successors(n)}, \cmd{predecessors(n)}; \cmd{n in graph} tests membership.

\api{ancestors(graph, source) -> set}
All nodes with a path to \cmd{source}, excluding it.

\api{is\_directed\_acyclic\_graph(graph) -> bool\quad find\_cycle(graph) -> list[(u, v)]}
Acyclicity test; one cycle as an edge list (\cmd{CycleError} if there is none).

\api{topological\_sort(graph, tie\_breaking\_priority=None) -> list\quad topological\_sort\_fifo(graph) -> list}
Nodes in topological order; \cmd{CycleError} on a cycle. Dependencies define only a partial
order; \cmd{tie\_breaking\_priority} maps nodes to comparable values (lower first, default
\cmd{()}) to make the order of independent nodes stable.

\api{class CycleError(Exception)}

% ======================================================================
\section{\texorpdfstring{\cmd{laupy.maxwell}}{laupy.maxwell}}
% ======================================================================

Categorises Maxwell nodes by hardware with the GPU/CPU model lists \cmd{GPU\_EXCELLENT},
\cmd{GPU\_STRONG}, \cmd{GPU\_WEAK}, \cmd{CPU\_EXCELLENT}, \cmd{CPU\_STRONG}, \cmd{CPU\_WEAK}
(\cmd{GPU\_ALL}, \cmd{CPU\_ALL}: all known models). Used by \cmd{submitslurm -g/-e/-s/-x} and
\cmd{listnodes}.

\textbf{Where the nodes come from.} SLURM tells which nodes exist and their states; the hardware
(CPU/GPU models, hence the classes) comes from the feature provider. The SLURM node features
(\cmd{sinfo -a -N} with \cmd{\%f}, \cmd{get\_live\_node\_table()}) take seconds to query, the node
states (\cmd{sinfo -N -o "\%N \%t"}) a fraction of a second, so by default
(\cmd{"app-managed-list"}) the hardware comes from the app-managed node table
\cmd{default\_node\_table\_path()} = \cmd{\$LAUPY\_MAXWELL\_TSV} or \cmd{\textasciitilde/.laupy/maxwell.tsv}:
unpacked from \cmd{PACKAGED\_NODE\_TABLE} (\cmd{laupy/maxwell.tsv.gz}) when missing and at once
updated from SLURM (retried on the next call if that fails), or created from SLURM without it; refreshed from SLURM when older than \cmd{MANAGED\_LIST\_MAX\_AGE\_DAYS} (7).
Messages about this are collected in \cmd{NODE\_NOTES}. When rows are combined, SLURM values win;
the table fills FLOPS, warranty and the hardware of nodes without features, resolves ambiguous
features and keeps nodes SLURM does not report. All queries are made at most once per process and
cached. A missing table is not an error; without SLURM and without a table the node lists are empty.

\api{set\_node\_table(path=None, live=True, use\_table=True, feature\_provider="app-managed-list")\quad find\_node\_table() -> str | None\quad get\_node\_source() -> dict}
Choose the sources, dropping the caches: \cmd{path} a table used as it is (\cmd{FileNotFoundError}
if missing; None: the managed table); \cmd{live=False} queries no SLURM (offline, no refresh);
\cmd{use\_table=False} no table at all (implies \cmd{"slurm"}); \cmd{feature\_provider}
\cmd{"app-managed-list"}, \cmd{"list"} (the table as it is, SLURM only for states) or \cmd{"slurm"}
(features on every call, completed from the table). A table provider without a table falls back
to \cmd{"slurm"}. \cmd{find\_node\_table()} returns the table to use (managing it if needed);
\cmd{get\_node\_source()} gives \cmd{\{"source", "table", "live\_error", "feature\_provider",
"hardware\_from"\}}. \cmd{node\_table\_candidates()}
lists the searched paths. Every SLURM query is timed in \cmd{QUERY\_LOG};
\cmd{format\_query\_log()} gives one line per query (shown by \cmd{submitslurm -{}-verbose}).

\api{get\_all\_nodes(partition=None)\quad get\_gpu\_nodes(partition=None)}
\api{get\_excellent\_nodes / get\_strong\_nodes / get\_weak\_nodes(partition=None)}
\api{get\_\{excellent,strong,weak\}\_\{gpu,cpu\}\_nodes(partition=None)}
Host names of each category, optionally limited to partitions (same forms as
\cmd{partition\_list()}). Built once per process (\cmd{init\_node\_cache()}) from the SLURM query,
or from the table when SLURM is switched off or fails.

\api{partition\_list(partition) -> list[str]}
Normalises a partition argument: \cmd{"allcpu,allgpu"} or a list/tuple/set of names; names are
stripped, empty ones dropped; \cmd{None} or \cmd{""} gives \cmd{[]}, meaning all partitions.
\raises \cmd{ValueError} for any other type.

\api{select\_nodes(partition=None, node\_class=None, gpu\_only=False, nodelist=None, state="any") -> list[str]}
Nodes of the partitions that SLURM reports, sorted: \cmd{node\_class} \cmd{"excellent"},
\cmd{"strong"} (includes excellent) or \cmd{"weak"}; \cmd{gpu\_only}; \cmd{nodelist} (SLURM syntax
or a list); \cmd{state} \cmd{"any"}, \cmd{"up"} (\cmd{node\_is\_up()}) or \cmd{"idle"}. One cached
query for classes and states. Used by \cmd{listnodes}.

\api{sbatch\_node\_request(partition=None, nodelist=None, gpu\_only=False, excellent=False, strong=False, exclude\_weak=False) -> NodeRequest}
The node options of \cmd{submitslurm}: \cmd{NodeRequest} has \cmd{nodelist} (required nodes, only
nodes that are up), \cmd{exclude} (weak nodes), \cmd{error} (no available node for \cmd{-n}/\cmd{-g}:
do not submit), \cmd{notes} (skipped nodes, fallback without restriction) and
\cmd{sbatch\_arguments()} (\cmd{-{}-nodelist=}/\cmd{-{}-exclude=}, never empty).

\api{get\_live\_nodes(partition=None, idle\_only=False) -> list[str]}
Nodes reported by \cmd{sinfo} in any state (optionally only idle ones), for the partitions given
as in \cmd{partition\_list()}. Reuses the query of the node lists, so \cmd{submitslurm} with node
options makes a single \cmd{sinfo} call.

\api{nodes\_from\_string(s) -> list[str]\quad expand\_node\_ranges(entry) -> list[str]}
Expand SLURM node lists: \cmd{"max-node[001-003,007]"} $\rightarrow$ individual host names (numbers zero-padded to three digits).

\api{compress\_node\_list(nodes) -> str}
The inverse: \cmd{["max-wng004", "max-wng005", "max-wng007"]} $\rightarrow$
\cmd{"max-wng[004-005,007]"}. Names without a three-digit suffix are kept as they are.

\api{get\_cluster\_topology(partition=None, live=True, table\_path=None) -> dict}
Describes the cluster (\cmd{listnodes -{}-topology}). By default SLURM is the source of truth
(\cmd{get\_live\_node\_table()}, the same query as the node lists); the node table, if any, only
fills gaps, resolves ambiguous features and is compared with it. If \cmd{sinfo} fails it falls back
to the table without states. \cmd{live=False} queries nothing and describes the table;
\cmd{table\_path} is a shortcut for \cmd{set\_node\_table(table\_path, live=live)}.
Keys:
\begin{description}[style=nextline,leftmargin=1.5em,itemsep=0pt]
  \item[\cmd{nodes}] One dict per described node in the partitions: Host, CPU, GPU (list), Memory,
    CPUCount, Partitions, Warranty, \cmd{class} (excellent, strong or weak; \emph{default} for hardware in \cmd{GPU\_ALL}/\cmd{CPU\_ALL}
    but in none of these classes; \emph{unclassified} for hardware not in \cmd{GPU\_ALL}/\cmd{CPU\_ALL}), \cmd{state}
    (compact \cmd{sinfo} state such as \cmd{idle}, \cmd{mix}, \cmd{drain}, \cmd{down*}) and
    \cmd{up}; the last two are \cmd{None} without live data.
  \item[\cmd{gpu\_types}] \cmd{\{type: \{"nodes": [...], "gpus": n\}\}} for nodes with GPUs.
  \item[\cmd{cpu\_types}] \cmd{\{cpu: [nodes]\}} for nodes without GPUs.
  \item[\cmd{classes}] Nodes per class, as used by \cmd{submitslurm -e/-s/-x}.
  \item[\cmd{not\_described}] Nodes \cmd{sinfo} reports that are missing in the table.
  \item[\cmd{table}] Path of the node table, \cmd{None} without one.
  \item[\cmd{unknown\_hardware}] \cmd{\{"GPU X" | "CPU Y": [nodes]\}}: hardware not in \cmd{GPU\_ALL}
    or \cmd{CPU\_ALL}, with the (unclassified) nodes that have it. Nodes with GPUs are judged by their GPUs,
    nodes without GPU by their CPU (\cmd{unknown\_hardware(node)}, \cmd{node\_class(node)}).
  \item[\cmd{source}, \cmd{partitions}, \cmd{live\_error}] \cmd{"sinfo"} or the path of the table used,
    the partitions used, and why \cmd{sinfo} could not be queried (\cmd{None} if it could).
    A failing \cmd{sinfo} does not raise.
  \item[\cmd{not\_in\_sinfo}, \cmd{hardware\_differs}, \cmd{filled\_from\_table}] SLURM source only:
    table nodes SLURM does not report; \cmd{(host, table hardware, SLURM hardware)} tuples; and
    \cmd{\{host: [fields]\}} taken from the table because SLURM gave no features.
\end{description}
\cmd{not\_described}, \cmd{not\_in\_sinfo}, \cmd{hardware\_differs} and \cmd{unknown\_hardware} show
where the node table or the model lists need updating.

\api{get\_live\_node\_table(table=None) -> list[dict]}
All nodes from \cmd{sinfo -a -N -o "\%N|\%P|\%t|\%c|\%z|\%m|\%G|\%f"} in the row format of
\cmd{parse\_maxwell\_tsv()}, plus \cmd{state} and \cmd{source} (fields taken from the table).
Partitions, state, CPUs (\cmd{"40 (2/10/2)"}), memory, GPU count (gres, else the \cmd{GPUxN}
feature) and the CPU/GPU models (features) come from SLURM; FLOPS, Warranty and the hardware of
nodes without features from \cmd{table} (default: the node table found, if any). CPU load
(\cmd{\%C}) is in \cmd{cpu\_load}.
\raises \cmd{RuntimeError} when \cmd{sinfo} is missing or fails.

\api{hardware\_from\_features(features, table\_row=None) -> (cpu, gpu\_model, gpu\_count)}
CPU and GPU model from node features, named as in \cmd{maxwell.tsv} and the \cmd{GPU\_*}/\cmd{CPU\_*}
lists: \cmd{Gold-6240} $\rightarrow$ \cmd{"Gold 6240 CPU @ 2.60GHz"}, \cmd{EPYC, 9374F} $\rightarrow$
\cmd{"AMD EPYC 9374"}, \cmd{E5-2640} with \cmd{V4} $\rightarrow$ \cmd{"E5-2640 v4 @ 2.40GHz"}. GPU
features are mapped by \cmd{GPU\_FEATURE\_NAMES} (\cmd{A100} $\rightarrow$ \cmd{GA100}, \ldots);
features in \cmd{GPU\_AMBIGUOUS\_FEATURES} (\cmd{RTX6000}: RTX A6000 or GB202GL RTX PRO) take the
model from \cmd{table\_row}; unknown features (e.g.\ \cmd{L40S}) are kept as they are and show up as
unclassified. \cmd{None} where the features say nothing.

\api{update\_node\_table(file\_path=None) -> (path, count, live\_rows)\quad write\_node\_table(file\_path, rows) -> int}
Refresh a table from a SLURM feature query (default: the managed one; created if missing), combined
with its old content (\cmd{listnodes -{}-update-default-list}, \cmd{-{}-update-list FILE}); and write
rows atomically (temporary file and rename, so parallel runs never read half a file).
\raises \cmd{RuntimeError} when SLURM cannot be queried.

\api{gres\_gpu\_count(gres) -> int | None\quad export\_topology(file\_path, rows=None) -> int}
GPUs in a \cmd{\%G} string (\cmd{"gpu:a100:4(S:0-1)"} $\rightarrow$ 4); and writing the nodes SLURM
reports, completed from the node table in use (or \cmd{rows}), as a \cmd{maxwell.tsv} file
(\cmd{listnodes -{}-export-topology}). Table nodes SLURM does not report are kept as they are in the table.
\raises \cmd{RuntimeError} when SLURM cannot be queried.

\api{format\_cluster\_topology(topology, verbose=False, user=None, color=False) -> str}
The text report printed by \cmd{listnodes -{}-topology}; \cmd{verbose} adds one line per node.
Given the result of \cmd{get\_cluster\_workload()} it prints the \cmd{-{}-workload} report, with the
jobs of \cmd{user} (default \cmd{current\_user()}) marked, in green with \cmd{color=True}. GPU types
are split by GPUs per node.

\api{get\_cluster\_workload(partition=None, table\_path=None) -> dict}
\cmd{get\_cluster\_topology()} plus the current load (\cmd{listnodes -{}-workload}). Each node also
gets \cmd{idle} (up, idle state, no CPU allocated), \cmd{cpus\_idle} and \cmd{cpus\_total}. The extra
key \cmd{workload} holds \cmd{error} (\cmd{None} or why \cmd{sinfo}/\cmd{squeue} failed; it does not
raise), \cmd{running\_jobs}, \cmd{running\_users} (\cmd{\{user: \{"jobs", "nodes", "pending"\}\}}),
\cmd{pending\_jobs}, \cmd{pending\_reasons}, \cmd{pending\_partitions} and \cmd{node\_users}
(\cmd{\{node: \{user: running jobs\}\}}). Running jobs are counted on the described nodes whatever
their partition; pending jobs only in the given partitions.

\api{get\_queue\_jobs(partition=None) -> list[dict]\quad get\_node\_cpu\_load(partition=None) -> dict\quad array\_task\_count(job\_id) -> int}
The \cmd{squeue} and \cmd{sinfo} queries behind it: jobs with \cmd{job\_id}, \cmd{user},
\cmd{state}, \cmd{partition}, \cmd{nodes}, \cmd{cpus}, \cmd{reason}, \cmd{tasks}; per-node CPUs
\cmd{\{"alloc", "idle", "other", "total"\}}; and the number of tasks of a pending array
(\cmd{"123\_[1-100\%10]"} $\rightarrow$ 100). The first two raise \cmd{RuntimeError} when the
command fails.

\api{get\_live\_node\_states(partition=None) -> dict[str, str]\quad node\_is\_up(state) -> bool\quad node\_class(node) -> str}
Per-node compact states from \cmd{sinfo -N}; whether a state can run jobs (\cmd{idle}, \cmd{mix},
\cmd{alloc}, \ldots, not down, drained or not responding \cmd{*}); and the class of a table row.
\raises \cmd{RuntimeError} from \cmd{get\_live\_node\_states()} when \cmd{sinfo} is missing or fails.

\api{parse\_maxwell\_tsv(file\_path) -> list[dict]\quad get\_node\_description()\quad run\_command(command)}
Table parsing (GPU column \cmd{"4xRTX 2080"} becomes a list of four names, partitions a list,
wiki-escaped \cmd{\textbackslash\_} in partition names becomes \cmd{\_}), the
cached table, and a subprocess helper returning stdout.

% ======================================================================
\section{\texorpdfstring{\cmd{laupy.redis\_server}}{laupy.redis\_server}}
% ======================================================================

Lifecycle of a temporary Redis server, as used for the GPU pool. The \cmd{REDIS} command is its
command line, and \cmd{GPU -{}-try-start-server} calls it directly, so neither depends on the other
being on \cmd{PATH}. Messages go through an \cmd{out} function (\cmd{print} by default; \cmd{GPU}
passes one writing to stderr, as its stdout carries only results). Problems are logged through the
module's logger \cmd{laupy.redis\_server}; the module configures no handler, the calling command
decides where log messages go. Failures raise \cmd{RedisLifecycleError} (a \cmd{RuntimeError}).

\api{RedisManager(host, port, redis\_server, pidfile, log\_dir, no\_log, out)}
Defaults: \cmd{host="127.0.0.1"}, \cmd{port=6379}, \cmd{redis\_server="redis-server"},
\cmd{pidfile=None}, \cmd{log\_dir="log"}, \cmd{no\_log=False}, \cmd{out=print}.
A Redis instance on \cmd{host:port}, identified by its PID file (default
\cmd{/tmp/laupy-gpu-<uid>-<SLURM job or local>-<port>.pid}). The log is timestamped
(\cmd{<log\_dir>/laupy-gpu-<time>-<job>-<port>.log}, \cmd{log\_dir} relative to the working
directory); \cmd{no\_log} writes none. \cmd{redis\_server} is a name searched in \cmd{PATH} or a path.

\api{start(kill\_unresponsive=False, extra\_config=(), dry\_run=False)}
Starts the server unless a responsive Redis already runs on the port (whoever started it). Persistence
is off and the database file name is timestamped, so no old state is loaded. A port held by your own
\cmd{redis-server} that does not respond is freed only with \cmd{kill\_unresponsive} (SIGTERM, then
SIGKILL); a port held by something else is refused. \cmd{dry\_run} only reports the configuration and
the files a start would create.

\api{stop(force=False, dump=False)\quad restart(...)\quad status()}
\cmd{stop} shuts down the server of the PID file (\cmd{SHUTDOWN NOSAVE}); \cmd{force} ignores the PID
file and terminates your own \cmd{redis-server} on the port; \cmd{dump} first saves the data to the
timestamped RDB file in \cmd{log\_dir}. \cmd{restart} takes the options of both. \cmd{status} reports
the port, the server, the PID file and the log, and logs inconsistencies as warnings.

\api{probe\_client(host, port, **client\_options) -> redis.Redis\quad ping(host, port) -> bool}
A client for quick checks of a server that may be hung: connecting gives up after
\cmd{PROBE\_CONNECT\_TIMEOUT} (1~s), each reply after \cmd{PROBE\_TIMEOUT} (3~s), with no retries
(redis-py $\geq$ 6 retries timed-out commands by default, about 38~s per call). Not for clients that
wait for work, such as \cmd{BRPOP}. \cmd{client\_options} are further \cmd{redis.Redis} arguments
(e.g.\ \cmd{db}, \cmd{decode\_responses}).

\api{describe\_port(host, port, verbose=False, out=print) -> dict\quad port\_owners(port) -> dict}
Who listens on the port (\cmd{psutil}, no \cmd{ss} needed; other users' processes are not visible, a
connection test tells such a port from a free one) and whether a responsive Redis answers:
\cmd{port\_state} \cmd{"free"}, \cmd{"redis"} or \cmd{"occupied"}, plus \cmd{owner\_unknown}.

% ======================================================================
\section{\texorpdfstring{\cmd{laupy.util}}{laupy.util}}
% ======================================================================

\api{list\_processes(pattern="kct-", skip\_pid=None) -> list[dict]}
Processes whose name matches \cmd{pattern}: \cmd{pid}, \cmd{cmdline}, \cmd{cwd}.

\api{get\_process\_info(pid) -> dict\quad print\_process\_info(proc\_info)}
\cmd{pid}, \cmd{ppid}, \cmd{name}, \cmd{cmdline}, \cmd{cwd} of one process; printing helper.

\appendix

% ======================================================================
\section{Entry fields}\label{sec:fields}
% ======================================================================

\begin{longtable}{@{}>{\raggedright}p{0.32\linewidth}>{\raggedright\arraybackslash}p{0.64\linewidth}@{}}
\toprule
Field & Meaning \\
\midrule
\endhead
\cmd{job\_id} & laupy entry ID, 32 hex characters (\cmd{uuid4().hex}). \\
\cmd{slurm\_id} & SLURM job ID (\cmd{int}). \\
\cmd{script\_name} & Script file name, e.g.\ \cmd{02destar.sbatch}. \\
\cmd{job\_name} & SLURM job name, \cmd{<script>\_<subdir>} in \cmd{submitslurm}. \\
\cmd{dag\_node\_name} & Node as string: \cmd{"3"}, \cmd{"reco"}, \cmd{"-1"}. \\
\cmd{step} & Step number; $-2$ for named nodes. \\
\cmd{dag\_node\_deps} & Parents of a named node (\cmd{[]} for numeric nodes). \\
\cmd{dependency\_slurm\_ids} & SLURM IDs the job waited for (\cmd{afterok}). \\
\cmd{command} & Script command, space-joined. \\
\cmd{slurm\_command} & Full \cmd{sbatch} argument list, used by \cmd{resubmit\_slurm\_job()}. \\
\cmd{laupy\_python} & Python environment (\cmd{LAUPY\_PYTHON}) exported to the job; absent for the default. \\
\cmd{processing\_unit\_subdir} & Absolute path of the processing unit. \\
\cmd{processing\_unit\_name} & Its directory name. \\
\cmd{timestamp} & Submission time, Unix seconds. \\
\cmd{slurm\_info} & Last known SLURM data, keys as from \cmd{slurm\_info()}. \\
\cmd{retired} & \cmd{True} once retired. \\
\cmd{retire\_timestamp} & When it was retired. \\
\cmd{retired\_reason} & Free text, e.g.\ \emph{Resubmitted}, \emph{pipeline retire}. \\
\cmd{retired\_by} & \cmd{job\_id} of the entry that replaced it. \\
\bottomrule
\end{longtable}

Fields written by laupy before 1.2.5: \cmd{execution\_unit\_dir} and \cmd{execution\_unit\_name}
are renamed on load.

% ======================================================================
\section{Behaviour worth knowing}
% ======================================================================

\begin{itemize}
  \item \textbf{States are only as fresh as the last refresh.} Selection, dependency resolution and
    cancellation use \cmd{slurm\_info}; refresh before deciding.
  \item \textbf{Final states are cached.} With \cmd{filter\_terminal\_states=True} (the default),
    entries in \cmd{TERMINAL\_STATES} are never queried again. Refresh with
    \cmd{filter\_terminal\_states=False} to re-read them.
  \item \textbf{Nothing saves implicitly} except \cmd{retire\_selected()}, \cmd{cancel\_selected()}
    and the \cmd{PipelineEntriesStore.store()} you call.
  \item \textbf{\cmd{schedule\_slurm\_job()} and \cmd{resubmit\_slurm\_job()} do not append} the new
    entry; append it, then store.
  \item \textbf{\cmd{sbatch} runs in the project root}, two levels above the processing unit, so
    relative paths in \cmd{slurm\_cmd} are relative to the root.
  \item \textbf{Logging:} laupy logs through \cmd{logging.getLogger("laupy.\ldots")}; configure the
    root logger to see the details of failed submissions.
\end{itemize}

\end{document}
